\documentclass[
  coursetitle={Cómputo Nube y Tecnologías Emergentes},
  coursehours={96},
  coursedate={Verano 2026},
  doctype={Taller},
  otherlangs={english,french},
  authorname={Lázaro},
  authorsurname={Bustio-Martínez},
  authortitle={PhD},
  role={Profesor Investigador},
  email={lbustio@inaoe.mx},
  phone={5559504000},
  ext={7459},
  institution={Instituto Nacional de Astrofísica, Óptica y Electrónica (INAOE)},
  department={Maestría en Ciencias en Tecnologías de Seguridad},
  address={Luis Enrique Erro 1, Sta. Ma. Tonantzintla, San Andrés Cholula, CP 72840, Puebla, México},
  margin={0.5in}
]{../../lbustio_lecture_notes}

\usepackage[utf8]{inputenc}
\usepackage[T1]{fontenc}
\usepackage{array}
\usepackage{tabularx}
\usepackage{enumitem}
\usepackage{longtable}
\usepackage{booktabs}
\usepackage{url}
\usepackage{listings}
\usepackage{xcolor}
\usepackage{graphicx}
\usepackage{tcolorbox}
\usepackage{verbatim}
\usepackage{tikz}
\usetikzlibrary{arrows.meta,backgrounds,fit,positioning}

\definecolor{lightgray}{rgb}{0.97, 0.97, 0.97}
\definecolor{darkblue}{rgb}{0, 0, 0.6}
\definecolor{gcpblue}{HTML}{4285F4}
\definecolor{gcpgreen}{HTML}{34A853}
\definecolor{gcpyellow}{HTML}{FBBC04}
\definecolor{gcpred}{HTML}{EA4335}
\definecolor{softgray}{HTML}{F5F7FA}
\definecolor{androidgreen}{HTML}{3DDC84}
\definecolor{androiddark}{HTML}{073042}

\lstset{
  backgroundcolor=\color{lightgray},
  basicstyle=\ttfamily\small,
  breaklines=true,
  frame=single,
  columns=fullflexible,
  showstringspaces=false,
  commentstyle=\color{gray},
  keywordstyle=\color{darkblue},
  inputencoding=utf8,
  extendedchars=true,
  literate={á}{{\'a}}1 {é}{{\'e}}1 {í}{{\'i}}1 {ó}{{\'o}}1 {ú}{{\'u}}1 {ñ}{{\~n}}1
           {Á}{{\'A}}1 {É}{{\'E}}1 {Í}{{\'I}}1 {Ó}{{\'O}}1 {Ú}{{\'U}}1 {Ñ}{{\~N}}1
}

\setlist[itemize]{topsep=0.3em,itemsep=0.2em}
\setlist[enumerate]{topsep=0.3em,itemsep=0.2em}

\begin{document}

\IberoCover

\section{Descripción de la sesión}

Esta sesión corresponde al primer taller autónomo de la práctica \textit{Detección de APKs maliciosas mediante permisos de Android}, introducida en la sesión del 15 de junio. El objetivo de hoy es avanzar de forma independiente en la primera mitad del análisis: preparar el entorno, cargar y limpiar el dataset, ejecutar el análisis exploratorio y aplicar agrupamiento no supervisado con K-Means.

Al finalizar esta sesión, el estudiante habrá completado las siguientes etapas de la práctica:

\begin{itemize}
    \item configuración del entorno de trabajo;
    \item carga y validación inicial del dataset;
    \item limpieza y construcción de la matriz de permisos;
    \item análisis exploratorio de datos (EDA);
    \item análisis de permisos sensibles y solapamiento entre clases;
    \item agrupamiento no supervisado con K-Means y visualización con PCA y t-SNE;
    \item contraste preliminar de las hipótesis H1, H2 y H3.
\end{itemize}

\begin{tcolorbox}[title=Referencia]
Toda la implementación se realiza en el notebook \texttt{notebooks/deteccion\_apks\_maliciosas\_permisos.ipynb} del proyecto \texttt{malware\_apk\_detection}, disponible desde la sesión del 15 de junio. Esta guía describe los objetivos de cada etapa y los resultados esperados; el código completo se encuentra en el notebook.
\end{tcolorbox}

\section{Preparación del entorno}

Active el entorno conda del proyecto y abra el notebook:

\begin{lstlisting}
conda activate apk-malware
jupyter lab notebooks/deteccion_apks_maliciosas_permisos.ipynb
\end{lstlisting}

Verifique que todas las dependencias estén disponibles ejecutando la primera celda de importaciones del notebook. Si algún paquete falta, instálelo con:

\begin{lstlisting}
pip install <nombre_del_paquete>
\end{lstlisting}

\begin{tcolorbox}[title=Punto de control 1]
\textbf{Resultado esperado:} la primera celda del notebook se ejecuta sin errores y todas las bibliotecas quedan importadas correctamente.
\end{tcolorbox}

\section{Carga y validación inicial del dataset}

Ejecute la celda de configuración de rutas y la celda de carga del dataset. Verifique las dimensiones, los tipos de datos y la distribución de la variable objetivo:

\begin{itemize}
    \item ¿Cuántas APKs tiene el dataset?
    \item ¿Cuántos permisos candidatos hay?
    \item ¿Existen valores faltantes o duplicados?
    \item ¿Cuántas APKs son benignas y cuántas maliciosas?
\end{itemize}

\begin{tcolorbox}[title=Punto de control 2]
\textbf{Resultado esperado:} 398 APKs, 330 permisos candidatos, dataset balanceado (199 benignas / 199 maliciosas), sin valores faltantes.
\end{tcolorbox}

\section{Limpieza y construcción de la matriz de permisos}

Ejecute las celdas de limpieza: eliminación de duplicados exactos, manejo de valores faltantes y conversión a tipo entero binario.

Registre en su bitácora:

\begin{itemize}
    \item ¿Cuántos duplicados se encontraron?
    \item ¿Cuántas APKs quedan tras la limpieza?
    \item ¿Cambió la distribución de clases? ¿Por qué?
\end{itemize}

\begin{tcolorbox}[title=Punto de control 3]
\textbf{Resultado esperado:} 224 APKs tras eliminar 174 duplicados exactos; 114 benignas y 110 maliciosas.
\end{tcolorbox}

\section{Análisis exploratorio de datos}

Ejecute las celdas de análisis exploratorio en el orden que aparecen en el notebook:

\subsection{Distribución de clases}

Genere el gráfico de barras con la distribución de clases tras la limpieza. Verifique que ambas clases tienen una representación equilibrada.

\subsection{Número de permisos por APK}

Ejecute las celdas que calculan y visualizan el número de permisos solicitados por clase. Registre:

\begin{itemize}
    \item media y mediana de permisos en APKs benignas;
    \item media y mediana de permisos en APKs maliciosas;
    \item observación visual del boxplot e histograma.
\end{itemize}

\begin{tcolorbox}[title=Evidencia para H2]
Si las APKs maliciosas solicitan en promedio más permisos que las benignas, esta diferencia constituye evidencia preliminar a favor de H2. Anote el valor exacto de ambas medias en su bitácora.
\end{tcolorbox}

\subsection{Permisos más frecuentes}

Identifique los 20 permisos más declarados en el dataset. Reflexione sobre cuáles de ellos podrían asociarse con comportamientos maliciosos.

\subsection{Permisos diferenciales entre clases}

Calcule y analice la diferencia de frecuencia de cada permiso entre clases. Liste los cinco permisos con mayor diferencia positiva (más asociados a malware) y reflexione sobre las capacidades del dispositivo que habilitan.

\section{Análisis de permisos sensibles}

Ejecute las celdas de análisis de permisos sensibles. Registre:

\begin{itemize}
    \item número promedio de permisos sensibles en APKs benignas;
    \item número promedio de permisos sensibles en APKs maliciosas;
    \item los tres permisos sensibles más frecuentes en malware.
\end{itemize}

\begin{tcolorbox}[title=Evidencia para H3]
Si las APKs maliciosas tienen significativamente más permisos sensibles en promedio, esto respalda H3. Compare sus resultados con el valor esperado del notebook: aproximadamente 9 permisos sensibles en maliciosas frente a 4 en benignas.
\end{tcolorbox}

\section{Solapamiento entre clases}

Ejecute las celdas que calculan permisos exclusivos de cada clase y permisos compartidos. Registre:

\begin{itemize}
    \item número de permisos exclusivos de APKs benignas;
    \item número de permisos exclusivos de APKs maliciosas;
    \item número de permisos compartidos entre ambas clases.
\end{itemize}

Reflexione sobre las implicaciones del solapamiento para la clasificación automática. ¿Por qué un permiso compartido no es evidencia directa de malware?

\section{Agrupamiento no supervisado con K-Means}

Ejecute las celdas de agrupamiento K-Means con dos clusters. Analice la tabla de contingencia (cluster vs. clase real):

\begin{itemize}
    \item ¿Un cluster concentra mayoritariamente APKs maliciosas?
    \item ¿El otro concentra mayoritariamente APKs benignas?
    \item ¿Qué tan puro es cada cluster?
\end{itemize}

Ejecute también las visualizaciones PCA y t-SNE. Observe si las dos clases forman regiones visualmente separables en las proyecciones bidimensionales.

\begin{tcolorbox}[title=Punto de control 4]
\textbf{Resultado esperado:} un índice Rand ajustado (ARI) alrededor de 0.45, indicando correspondencia moderada entre clusters y clases. Las proyecciones PCA y t-SNE deben mostrar cierta separación entre clases, aunque con solapamiento visible.
\end{tcolorbox}

\section{Contraste preliminar de hipótesis}

Con base en los resultados obtenidos hasta este punto, complete el contraste parcial de hipótesis en su bitácora:

\begin{description}
    \item[H2.] ¿El p-value de la prueba Mann-Whitney permite rechazar la hipótesis nula? ¿Cuál fue el p-value obtenido?
    \item[H3.] ¿Las APKs maliciosas tienen más permisos sensibles en promedio? ¿Cuáles son los tres más frecuentes?
    \item[H4 (parcial).] ¿Cuántos permisos son compartidos entre clases? ¿Qué implica este solapamiento?
\end{description}

\begin{tcolorbox}[title=Nota sobre H1]
H1 afirma que los modelos supervisados superan la línea base. Esta hipótesis no puede contrastarse todavía, ya que el entrenamiento de modelos se realizará en la siguiente sesión.
\end{tcolorbox}

\section{Entregables al cierre de esta sesión}

Al terminar la sesión, el estudiante debe tener completadas y documentadas las siguientes evidencias en su bitácora técnica:

\begin{enumerate}
    \item Dimensiones del dataset original y tras la limpieza.
    \item Gráficos del análisis exploratorio: distribución de clases, boxplot e histograma de permisos, top 20 permisos frecuentes, permisos diferenciales.
    \item Tabla comparativa de permisos sensibles por clase.
    \item Conteo de permisos exclusivos y compartidos entre clases.
    \item Tabla de contingencia del agrupamiento K-Means.
    \item Proyecciones PCA y t-SNE coloreadas por clase real.
    \item Contraste parcial de H2, H3 y H4 con evidencias numéricas.
    \item Registro de problemas técnicos encontrados y soluciones aplicadas.
\end{enumerate}

\begin{tcolorbox}[title=Avance hacia la sesión del 22/06]
En la próxima sesión se completará la segunda mitad de la práctica: entrenamiento de modelos supervisados, evaluación comparativa, contraste de las hipótesis H1, H5 y H6, prueba de la aplicación web con una APK real y elaboración del reporte final. Es indispensable llegar a esa sesión con todas las evidencias de esta sesión documentadas en la bitácora.
\end{tcolorbox}

\end{document}
